\chapter{Limitaciones, aspectos eticos y reproducibilidad}

\section{Funcion de este capitulo}

Este capitulo reune tres elementos que no son el centro del trabajo, pero si condicionan su interpretacion final: las limitaciones del estudio, el uso responsable del sistema y las condiciones para repetir el trabajo.

Algunas de estas restricciones ya se adelantan de forma breve en la metodologia
porque afectan a la interpretacion de la evaluacion. Aqui se recogen de forma
m\'as amplia como limites del sistema y del propio trabajo.

\section{Limitaciones del estudio}

\textit{Completar aqui limites de datos, dependencia del modelo, cobertura de casos, sensibilidad a prompts, dificultad de generalizacion y cualquier restriccion importante del sistema. Este es el lugar principal para reunir las limitaciones globales del trabajo y ampliar las restricciones que en metodologia solo se introducen para acotar la lectura de la evaluacion.}

\section{Aspectos eticos y uso responsable}

\textit{Completar aqui finalidad prevista, usos excluidos, supervision humana, limites del analisis historico y riesgos de ejecutar codigo generado o de sobrerinterpretar resultados financieros.}

\section{Reproducibilidad}

\textit{Completar aqui lo necesario para repetir el trabajo: entorno, dependencias, datos, artefactos conservados, variables de entorno, pasos minimos de ejecucion y factores que pueden hacer variar el resultado.}

\section{Preguntas de control}

\begin{enumerate}
    \item ?Que no deberia concluir un lector a partir del sistema?
    \item ?Que necesita exactamente otra persona para repetir un caso?
    \item ?Que parte del resultado depende del flujo y que parte depende del modelo concreto?
    \item ?Que riesgos siguen abiertos aunque el sistema funcione?
\end{enumerate}
